\subsection{Positive Semidefinite and Positive Definite Matrices}

\begin{definition} \label{def:positive-semidefinite-matrix} \hlt{Positive Semidefinite and Positive Definite Matrices}\\
Let $\mathbb{F}\in\{\mathbb{R},\mathbb{C}\}$ and let $A\in\operatorname{Mat}(n,n,\mathbb{F})$.\\
The matrix $A$ is \hlt{positive semidefinite}, written $A\succeq0$, if $A=A^*$ and $\mathbf{x}^*A\mathbf{x}\geq0$ for every $\mathbf{x}\in\mathbb{F}^n$.\\
It is \hlt{positive definite}, written $A\succ0$, if $A=A^*$ and $\mathbf{x}^*A\mathbf{x}>0$ for every nonzero $\mathbf{x}\in\mathbb{F}^n$.\\
Negative semidefinite and negative definite mean $-A\succeq0$ and $-A\succ0$, respectively.
\end{definition}

\begin{proposition} \label{prop:rectangular-gram-products-are-psd} \hlt{Rectangular Gram Products Are Positive Semidefinite}\\
Let $R\in\operatorname{Mat}(m,n,\mathbb{F})$, where $\mathbb{F}\in\{\mathbb{R},\mathbb{C}\}$. Then $R^*R\succeq0$, $RR^*\succeq0$, and $(R^*R)_{j,j}=\norm{R_{\cdot,j}}_2^2$; $(RR^*)_{i,i}=\norm{R_{i,\cdot}}_2^2$.\\
Full column rank gives $R^*R\succ0$; full row rank gives $RR^*\succ0$.
\begin{proof}
The matrices are self-adjoint by Proposition \ref{prop:basic-adjoint-identities}.\\
For $\mathbf{x}\in\mathbb{F}^n$ and $\mathbf{y}\in\mathbb{F}^m$, $\mathbf{x}^*R^*R\mathbf{x}=\norm{R\mathbf{x}}_2^2\geq0$ and $\mathbf{y}^*RR^*\mathbf{y}=\norm{R^*\mathbf{y}}_2^2\geq0$. Hence both products are positive semidefinite. Taking $\mathbf{x}=\mathbf{e}_j$ and $\mathbf{y}=\mathbf{e}_i$ gives the diagonal formulas. If $R$ has full column rank, then $R\mathbf{x}\neq0$ for $\mathbf{x}\neq0$, so $R^*R\succ0$. If $R$ has full row rank, apply the same argument to $R^*$.\\
\hlt{Structural Commentary}: Multiplying a matrix by its adjoint produces a Gram matrix of columns or rows.\\
\hlt{Key Insight}: The quadratic form is a squared norm, so no cancellation can make it negative.
\end{proof}
\end{proposition}

\begin{theorem} \label{thm:positive-semidefinite-matrix-equivalences} \hlt{Equivalent Forms of Positive Semidefiniteness}\\
Let $A\in\operatorname{Mat}(n,n,\mathbb{F})$, where $\mathbb{F}\in\{\mathbb{R},\mathbb{C}\}$. The following are equivalent:
\begin{enumerate}[label=(\roman*)]
\setlength{\itemsep}{0pt}
\item $A\succeq0$.
\item $A$ is Hermitian and all eigenvalues of $A$ are nonnegative.
\item $A=B^*B$ for some $B\in\operatorname{Mat}(m,n,\mathbb{F})$.
\item $A$ has a positive semidefinite square root.
\end{enumerate}
In this case the positive semidefinite square root is unique.\\
If $A=U\Lambda U^*$ with $\Lambda=\operatorname{diag}(\lambda_1,\ldots,\lambda_n)$ and $\lambda_i\geq0$, then $A^{1/2}=U\operatorname{diag}(\sqrt{\lambda_1},\ldots,\sqrt{\lambda_n})U^*$.
\begin{proof}
\begin{enumerate}[label=(\roman*).]
\setlength{\itemsep}{0pt}
\item[$(i)$]$\Rightarrow (ii)$ This is the matrix form of Theorem \ref{thm:positive-operator-characterization}.
\item[$(ii)$]$\Rightarrow (iv)$ By the spectral theorem, write $A=U\Lambda U^*$, where $U$ is unitary and $\Lambda=\operatorname{diag}(\lambda_1,\ldots,\lambda_n)$ with $\lambda_i\geq0$. Then $A^{1/2}=U\operatorname{diag}(\sqrt{\lambda_1},\ldots,\sqrt{\lambda_n})U^*$ is positive semidefinite and satisfies $(A^{1/2})^2=A$.
\item[$(iv)$]$\Rightarrow (iii)$ If $C\succeq0$ and $C^2=A$, then $C=C^*$ and $A=C^*C$, so (iii) holds with $m=n$.
\item[$(iii)$]$\Rightarrow (i)$ If $A=B^*B$, then $A^*=A$ and $\mathbf{x}^*A\mathbf{x}=\mathbf{x}^*B^*B\mathbf{x}=\norm{B\mathbf{x}}_2^2\geq0$, so $A\succeq0$.
\end{enumerate}
The positive semidefinite square root is unique by Proposition \ref{prop:positive-square-root-unique}.\\
\hlt{Structural Commentary}: Positive semidefinite matrices are Gram-type matrices with nonnegative spectrum.\\
\hlt{Key Insight}: Positivity is best checked either on eigenvalues or by finding a factor $B$ with $A=B^*B$.
\end{proof}
\end{theorem}

\begin{theorem} \label{thm:gram-factor-kernel-rank-column-data} \hlt{Kernel, Rank, and Columns of $B^*B$}\\
Let $B\in\operatorname{Mat}(m,n,\mathbb{F})$, where $\mathbb{F}\in\{\mathbb{R},\mathbb{C}\}$, and set $A=B^*B$. Then:
\begin{enumerate}[label=(\roman*)]
\setlength{\itemsep}{0pt}
\item $A\succeq0$;
\item $\ker(A)=\ker(B)$;
\item $\operatorname{rank}(A)=\operatorname{rank}(B)$;
\item $\operatorname{col}(A)=\operatorname{col}(B^*)$;
\item $A=0$ if and only if $B=0$.
\end{enumerate}
\begin{proof}
\begin{enumerate}[label=(\roman*)]
\setlength{\itemsep}{0pt}
\item Positive semidefiniteness follows from Theorem \ref{thm:positive-semidefinite-matrix-equivalences}.
\item If $B\mathbf{x}=0$, then $A\mathbf{x}=0$. Conversely, if $A\mathbf{x}=0$, then $0=\mathbf{x}^*A\mathbf{x}=\norm{B\mathbf{x}}_2^2$ so $B\mathbf{x}=0$.
\item Since $A$ and $B$ both have $n$ columns, rank-nullity and (ii) give $\operatorname{rank}(A)=\operatorname{rank}(B)$.
\item Since $A=A^*$, Proposition \ref{prop:adjoint-kernel-range-orthogonality} and (ii) give $\operatorname{col}(A)=\ker(A)^{\perp}=\ker(B)^{\perp}=\operatorname{col}(B^*)$.
\item By (iii), $\operatorname{rank}(A)=0$ if and only if $\operatorname{rank}(B)=0$.\\
A matrix has rank zero exactly when it is the zero matrix, so $A=0$ if and only if $B=0$.
\end{enumerate}
\hlt{Structural Commentary}: The Gram product $B^*B$ preserves the kernel of $B$.\\
\hlt{Key Insight}: The identity $\mathbf{x}^*B^*B\mathbf{x}=\norm{B\mathbf{x}}_2^2$ is the whole mechanism.
\end{proof}
\end{theorem}

\begin{figure}[H]
  \centering
  \begin{tikzpicture}[
  >=Stealth,
  every node/.style={font=\small},
  box/.style={draw=gray!55, rounded corners=3pt, fill=gray!6, align=center, inner xsep=6pt, inner ysep=4.5pt, minimum height=0.82cm},
  result/.style={box, fill=green!7, minimum width=2.4cm},
  arrow/.style={->, line width=0.85pt, blue!70!black}
]
  \node[box, minimum width=2.05cm] (factor) at (0,0) {$A=B^*B$};
  \node[box, minimum width=4.0cm] (test) at (3.55,0)
    {$\mathbf{x}^*A\mathbf{x}=\norm{B\mathbf{x}}_2^2$\\[-2pt]
     $\norm{B\mathbf{x}}_2^2=0\Longleftrightarrow B\mathbf{x}=0$};
  \node[result] (kernel) at (7.2,0) {$\ker(A)=\ker(B)$};

  \draw[arrow] (factor) -- (test);
  \draw[arrow] (test) -- (kernel);
\end{tikzpicture}

  \caption{For $A=B^*B$, the squared-norm identity gives $\ker A=\ker B$.}
  \label{fig:psd-factorization-kernel}
\end{figure}

\begin{corollary} \label{cor:positive-semidefinite-zero-quadratic-form} \hlt{Zero Quadratic Form Detects the Kernel}\\
If $A\in\operatorname{Mat}(n,n,\mathbb{F})$ is positive semidefinite, where $\mathbb{F}\in\{\mathbb{R},\mathbb{C}\}$, then $\mathbf{x}^*A\mathbf{x}=0$ if and only if $A\mathbf{x}=0$.
\begin{proof}
Write $A=B^*B$.
\begin{enumerate}[label=\roman*.]
\setlength{\itemsep}{0pt}
\item[$(\Rightarrow)$] If $\mathbf{x}^*A\mathbf{x}=0$, then $\norm{B\mathbf{x}}_2^2=0$, so $B\mathbf{x}=0$ and hence $A\mathbf{x}=B^*B\mathbf{x}=0$.
\item[$(\Leftarrow)$] If $A\mathbf{x}=0$, then $\mathbf{x}^*A\mathbf{x}=0$.
\end{enumerate}
\hlt{Structural Commentary}: A zero PSD quadratic value has no hidden cancellation.\\
\hlt{Key Insight}: Zero length-square means the factored vector itself is zero.
\end{proof}
\end{corollary}

\begin{method} \label{meth:psd-spectral-factorization} \hlt{PSD Spectral Factorisation Method}\\
To analyse a Hermitian matrix $A$:
\begin{enumerate}[label=(\roman*)]
\setlength{\itemsep}{0pt}
\item diagonalise $A=U\Lambda U^*$;
\item check the signs of the diagonal entries of $\Lambda$;
\item if $\Lambda\geq0$, set $A^{1/2}=U\Lambda^{1/2}U^*$ and use $A=(A^{1/2})^*A^{1/2}$;
\item translate kernel, rank, and zero-quadratic-form questions into the factor $A^{1/2}$.
\end{enumerate}
\end{method}

\begin{theorem} \label{thm:positive-definite-matrix-equivalences} \hlt{Equivalent Forms of Positive Definiteness}\\
Let $A\in\operatorname{Mat}(n,n,\mathbb{F})$, where $\mathbb{F}\in\{\mathbb{R},\mathbb{C}\}$. The following are equivalent:
\begin{enumerate}[label=(\roman*)]
\setlength{\itemsep}{0pt}
\item $A\succ0$.
\item $A$ is Hermitian and all eigenvalues of $A$ are positive.
\item $A\succeq0$ and $A$ is invertible.
\item $A=B^*B$ for some invertible $B\in\operatorname{Mat}(n,n,\mathbb{F})$.
\item $A$ is invertible and $A^{-1}\succ0$.
\end{enumerate}
\begin{proof}
\begin{enumerate}[label=(\roman*).]
\setlength{\itemsep}{0pt}
\item[$(i)$]$\Rightarrow (ii)$ By definition $A$ is Hermitian. If $(\lambda,\mathbf{v})$ is an eigenpair with $\mathbf{v}\neq\mathbf{0}$, then $\lambda\norm{\mathbf{v}}_2^2=\mathbf{v}^*A\mathbf{v}>0$, so every eigenvalue is positive.
\item[$(ii)$]$\Rightarrow (iii)$ Positive eigenvalues imply $A\succeq0$ by Theorem \ref{thm:positive-semidefinite-matrix-equivalences}; they also exclude $0$, so $A$ is invertible.
\item[$(iii)$]$\Rightarrow (iv)$ By Theorem \ref{thm:positive-semidefinite-matrix-equivalences}, $A$ has a positive semidefinite square root $A^{1/2}$.\\
Since $A$ is invertible, $A^{1/2}$ is invertible, and $A=(A^{1/2})^*A^{1/2}$.
\item[$(iv)$]$\Rightarrow (i)$ If $A=B^*B$ and $B$ is invertible, then for every nonzero $\mathbf{x}$, $\mathbf{x}^*A\mathbf{x}=\norm{B\mathbf{x}}_2^2>0$.
\item[$(ii)$]$\Leftrightarrow (v)$ Hermitian-ness is preserved by inversion, and eigenvalues of $A^{-1}$ are reciprocal eigenvalues of $A$. Thus positive eigenvalues of $A$ are equivalent to positive eigenvalues of $A^{-1}$.
\end{enumerate}
\hlt{Structural Commentary}: Positive definite matrices are the invertible part of the positive semidefinite cone.\\
\hlt{Key Insight}: Positive definite means every eigenvalue is strictly on the positive side of zero.
\end{proof}
\end{theorem}

\begin{proposition} \label{prop:trace-determinant-positive-semidefinite} \hlt{Trace and Determinant of Positive Semidefinite Matrices}\\
Let $A\in\operatorname{Mat}(n,n,\mathbb{F})$ be positive semidefinite, where $\mathbb{F}\in\{\mathbb{R},\mathbb{C}\}$. Then $\operatorname{tr}A\geq0$, with equality if and only if $A=0$, and $\det A\geq0$, with strict inequality if and only if $A$ is positive definite.
\begin{proof}
List the eigenvalues of $A$ as $\lambda_1,\ldots,\lambda_n$. By Theorem \ref{thm:positive-semidefinite-matrix-equivalences}, each $\lambda_j\geq0$.\\
Hence $\operatorname{tr}A=\sum\limits_{j=1}^{n}\lambda_j\geq0$ and $\det A=\prod\limits_{j=1}^{n}\lambda_j\geq0$.\\
The trace is zero exactly when every $\lambda_j=0$, which for a Hermitian matrix means $A=0$. The determinant is positive exactly when every $\lambda_j>0$, which is positive definiteness by Theorem \ref{thm:positive-definite-matrix-equivalences}.\\
\hlt{Structural Commentary}: Trace and determinant aggregate the nonnegative eigenvalues by sum and product.\\
\hlt{Key Insight}: Equality is controlled by which eigenvalues vanish.
\end{proof}
\end{proposition}

\begin{theorem} \label{thm:positive-semidefinite-stability-rules} \hlt{Stability Rules for Positive Semidefinite Matrices}\\
Let $A,B\succeq0$ over $\mathbb{F}\in\{\mathbb{R},\mathbb{C}\}$. Whenever the expressions are defined:
\begin{enumerate}[label=(\roman*)]
\setlength{\itemsep}{0pt}
\item $aA+bB\succeq0$ for $a,b\geq0$; if $a,b>0$ and one of $A,B$ is positive definite, then $aA+bB\succ0$.
\item $S^*AS\succeq0$; moreover, $S^*AS\succ0$ if and only if $S\mathbf{x}\notin\ker(A)$ for all $\mathbf{x}\neq\mathbf{0}$.
\item $A\oplus B\succeq0$, with $A\oplus B\succ0$ if and only if $A\succ0$ and $B\succ0$.
\end{enumerate}
\begin{proof}
\begin{enumerate}[label=(\roman*)]
\setlength{\itemsep}{0pt}
\item $\mathbf{x}^*(aA+bB)\mathbf{x}=a\mathbf{x}^*A\mathbf{x}+b\mathbf{x}^*B\mathbf{x}\geq0$. If, for example, $A\succ0$ and $a>0$, then for every nonzero $\mathbf{x}$ the first summand is positive and the second is nonnegative, so $aA+bB\succ0$; the case $B\succ0$ is identical.
\item $\mathbf{x}^*S^*AS\mathbf{x}=(S\mathbf{x})^*A(S\mathbf{x})\geq0$. By Corollary \ref{cor:positive-semidefinite-zero-quadratic-form}, equality occurs exactly when $S\mathbf{x}\in\ker(A)$.
\item $(\mathbf{x}\oplus\mathbf{y})^*(A\oplus B)(\mathbf{x}\oplus\mathbf{y})=\mathbf{x}^*A\mathbf{x}+\mathbf{y}^*B\mathbf{y}\geq0$. This expression is positive for every nonzero $\mathbf{x}\oplus\mathbf{y}$ exactly when each summand is positive on its own nonzero vectors, namely when $A\succ0$ and $B\succ0$.
\end{enumerate}
\hlt{Structural Commentary}: The PSD cone is stable under sums, pullbacks, and direct sums.\\
\hlt{Key Insight}: Test the quadratic form after each construction.
\end{proof}
\end{theorem}

\begin{proposition} \label{prop:two-by-two-positive-semidefinite-square-root} \hlt{Explicit Square Root for a $2\times2$ Positive Semidefinite Matrix}\\
Let $A\in\operatorname{Mat}(2,2,\mathbb{F})$, where $\mathbb{F}\in\{\mathbb{R},\mathbb{C}\}$, be positive semidefinite and nonzero. Set $\tau=\bigl(\operatorname{tr}A+2\sqrt{\det A}\bigr)^{1/2}$.\\
Then $\tau>0$ and $A^{1/2}=\frac{1}{\tau}\bigl(A+\sqrt{\det A}\,I_2\bigr)$.
\begin{proof}
Let $C=\frac{1}{\tau}\bigl(A+\sqrt{\det A}\,I_2\bigr)$. By Proposition \ref{prop:trace-determinant-positive-semidefinite}, $\det A\geq0$ and $\operatorname{tr}A>0$, hence $\tau>0$; by Theorem \ref{thm:positive-semidefinite-stability-rules}, $C\succeq0$. Cayley--Hamilton gives $A^2-(\operatorname{tr}A)A+(\det A)I_2=0$, so
\begin{align}
C^2 =\frac{1}{\tau^2}\bigl(A^2+2\sqrt{\det A}\,A+(\det A)I_2\bigr) =\frac{1}{\tau^2}\bigl((\operatorname{tr}A)A+2\sqrt{\det A}\,A\bigr)=A.
\nonumber
\end{align}
Thus $C=A^{1/2}$ by uniqueness of the positive semidefinite square root.\\
\hlt{Structural Commentary}: In dimension $2$, Cayley--Hamilton makes $A^{1/2}$ linear in $A$.\\
\hlt{Key Insight}: Square the linear candidate and use Cayley--Hamilton.
\end{proof}
\end{proposition}

\begin{definition} \label{def:directed-incidence-matrix} \hlt{Directed Incidence Matrix}\\
For edges $e=u\to v$ on vertices $1,\ldots,n$, $B\in\operatorname{Mat}(m,n,\mathbb{R})$ has row $e$ equal to $-\mathbf{e}_u^t+\mathbf{e}_v^t$.
\end{definition}

\begin{figure}[H]
  \centering
  \begin{tikzpicture}[
  >=Stealth,
  every node/.style={font=\small},
  vertex/.style={circle, draw=blue!65!black, fill=blue!7, minimum size=0.62cm},
  edge/.style={->, blue!70!black, line width=0.75pt},
  mat/.style={draw=gray!55, fill=gray!5, rounded corners=2pt, inner sep=5pt, align=center}
]
  \node[vertex] (v1) at (0,1.2) {$1$};
  \node[vertex] (v2) at (2.0,1.2) {$2$};
  \node[vertex] (v3) at (2.0,-0.6) {$3$};
  \node[vertex] (v4) at (0,-0.6) {$4$};

  \draw[edge] (v1) -- node[above] {$e_1$} (v2);
  \draw[edge] (v2) -- node[right] {$e_2$} (v3);
  \draw[edge] (v4) -- node[below] {$e_3$} (v3);
  \draw[edge] (v1) -- node[left] {$e_4$} (v4);

  \node[mat] at (5.0,0.3) {$B=\begin{bmatrix}
  -1&1&0&0\\
  0&-1&1&0\\
  0&0&1&-1\\
  -1&0&0&1
  \end{bmatrix}$};
\end{tikzpicture}

  \caption{A directed incidence matrix records oriented edge differences from vertices to edges.}
  \label{fig:directed-incidence-matrix}
\end{figure}

\begin{example} \label{ex:directed-incidence-matrix-four-node} \hlt{A Directed Incidence Matrix}\\
For Figure \ref{fig:directed-incidence-matrix}, $B=
\begin{bmatrix}
-1 & 1 & 0 & 0\\
0 & -1 & 1 & 0\\
0 & 0 & 1 & -1\\
-1 & 0 & 0 & 1
\end{bmatrix}$. For node potentials $\mathbf{x}=(x_1,x_2,x_3,x_4)^t$, $B\mathbf{x}=
\begin{bmatrix}
x_2-x_1\\
x_3-x_2\\
x_3-x_4\\
x_4-x_1
\end{bmatrix}$.\\
Hence $\ker(B)=\operatorname{span}\{(1,1,1,1)^t\}$, and $\operatorname{rank}(B)=3$. The loop vector $\mathbf{y}=(1,1,-1,-1)^t$ spans $\ker(B^t)$.\\
Thus $\mathbf{b}\in\operatorname{col}(B)$ if and only if $\mathbf{y}^t\mathbf{b}=b_1+b_2-b_3-b_4=0$. This is Kirchhoff loop compatibility.\\
\hlt{Structural Commentary}: Incidence subspaces encode graph constraints.\\
\hlt{Key Insight}: A sparse matrix can encode both the algebra of a linear map and the combinatorics of a graph.
\end{example}

\begin{perspective} \label{per:incidence-least-squares-ranking} \hlt{Incidence Data and Least Squares}\\
For an incidence matrix $B$ as in Example \ref{ex:directed-incidence-matrix-four-node}, $B\mathbf{x}=\mathbf{b}$ seeks node potentials matching edge differences. If loop consistency fails, least squares projects $\mathbf{b}$ onto $\operatorname{col}(B)$, where $B^*(B\mathbf{x}_0-\mathbf{b})=\mathbf{0}$ (Theorem \ref{thm:matrix-least-squares-normal-equations} on $B$). The constant-vector kernel means potentials are unique only after fixing one node or choosing minimum norm.
\end{perspective}

\begin{proposition} \label{prop:incidence-laplacian-positive-semidefinite} \hlt{Incidence Laplacians Are Positive Semidefinite}\\
Let $B$ be a directed incidence matrix and set $L=B^tB$. Then $L\succeq0$ and $\mathbf{x}^tL\mathbf{x}=\sum\limits_{e=(u\to v)}(x_v-x_u)^2$, with $\ker(L)$ equal to the vectors constant on each connected component.
\begin{proof}
Since $L=B^tB$, Proposition \ref{prop:rectangular-gram-products-are-psd} gives $L\succeq0$.\\
For $e=(u\to v)$, $(B\mathbf{x})_e=x_v-x_u$, so $\mathbf{x}^tL\mathbf{x}=\norm{B\mathbf{x}}_2^2=\sum\limits_{e=(u\to v)}(x_v-x_u)^2$. Thus $\mathbf{x}\in\ker(L)$ if and only if $B\mathbf{x}=\mathbf{0}$, equivalently $x_u=x_v$ on every edge, which means $\mathbf{x}$ is constant on each connected component.\\
\hlt{Structural Commentary}: Laplacians sum squared edge differences.\\
\hlt{Key Insight}: A graph Laplacian is a Gram product, so positivity is automatic.
\end{proof}
\end{proposition}
